\Askhsh{32225_SOLUTION}{1}
\begin{ylh}{1}
    \item [Ύλη από εικόνα]
\end{ylh}
\textbf{ΛΥΣΗ}
\begin{alist}
\item [\textalpha)]
\begin{rlist}
\item [i.] Για κάθε $x \in [-1,+\infty)$ είναι:
\[ (f(x)+x)^2 = x^2(x+1) \Leftrightarrow g^2(x)=x^2(x+1) \Leftrightarrow |g(x)|=|x|\sqrt{x+1}. \]
Είναι $g(x)=0 \Leftrightarrow |g(x)|=0 \Leftrightarrow |x|\sqrt{x+1}=0 \Leftrightarrow x=0 \text{ ή } x=-1$.
\item [ii.] Η συνεχής (ως άθροισμα συνεχών) $g(x)$ είναι μη μηδενιζόμενη σε καθένα από τα διαστήματα $(-1,0), (0,+\infty)$ άρα διατηρεί σταθερό πρόσημο σε καθένα από τα $(-1,0), (0,+\infty)$.
Ακόμη είναι $g(1)=f(1)+1 > 0$ και $g\PARENS{-\frac{1}{2}} = f\PARENS{-\frac{1}{2}} - \frac{1}{2} < 0$ αφού $f(1)>-1$ και $f\PARENS{-\frac{1}{2}} < \frac{1}{2}$.
Οπότε είναι $g(x)<0$ για κάθε $x \in [-1,0)$ και $g(x)>0$ για κάθε $x \in (0,+\infty)$.
\end{rlist}

\item [\textbeta)] Για $x \in (0,+\infty)$ είναι $|g(x)|=g(x)$ και $|x|=x$, έτσι έχουμε:
\[
g(x)=x\sqrt{x+1} \Leftrightarrow f(x)+x=x\sqrt{x+1}
\]
\[
\Leftrightarrow f(x)=x\sqrt{x+1}-x
\]
\[
\Leftrightarrow f(x)=x(\sqrt{x+1}-1)
\]
Για $x \in (-1,0)$ είναι $|g(x)|=-g(x)$ και $|x|=-x$, έτσι έχουμε:
\[
-g(x)=-x\sqrt{x+1} \Leftrightarrow g(x)=x\sqrt{x+1}
\]
\[
\Leftrightarrow f(x)=x(\sqrt{x+1}-1)
\]
Ακόμη είναι $f(0)=0 \cdot (\sqrt{0+1}-1)=0$ και $f(-1)=-1 \cdot (\sqrt{-1+1}-1)=1$ άρα ο τύπος της συνάρτησης $f$ είναι $f(x)=x(\sqrt{x+1}-1)$ για κάθε $x \ge -1$.

\item [\textgamma)] Η συνεχής συνάρτηση $f$ είναι κυρτή άρα η $f'$ είναι γνησίως αύξουσα στο $(-1,+\infty)$.
Ακόμη είναι $h'(x)=f'(x+1)\cdot(x+1)'-f'(x)=f'(x+1)-f'(x)$.
Για κάθε $x \in (-1,+\infty)$ είναι $x < x+1 \Rightarrow f'(x) < f'(x+1) \Rightarrow f'(x+1)-f'(x) > 0 \Rightarrow h'(x) > 0$ άρα η συνάρτηση $h$ είναι γνησίως αύξουσα.
Για κάθε $x \in [2023,2024]$ είναι $x < x+1$ άρα $h(x) < h(x+1)$.
Επιπλέον οι $h(x), h(x+1)$ είναι συνεχείς άρα ισχύει:
\[
\int_{2023}^{2024} h(x)dx < \int_{2023}^{2024} h(x+1)dx \Rightarrow \int_{2023}^{2024} (f(x+1)-f(x))dx < \int_{2023}^{2024} (f(x+2)-f(x+1))dx.
\]
\end{alist}